\documentclass[../../../include/open-logic-section]{subfiles}

\begin{document}

\olfileid[pt]{sth}{z}{infinity-again}	
\olsection{Infinito}

Já temos axiomas suficientes para garantir que há infinitamente muitos
conjuntos (se é que existe algum). De fato, suponhamos que algum conjunto existe; então
$\emptyset$ existe (por \olref[sth][z][sep]{prop:emptyexists}). Agora, para
qualquer conjunto~$x$, o conjunto $x \cup \{x\}$ existe por
\olref[sth][z][pairs]{prop:pairsconsequences}. Logo, aplicando isto algumas
vezes, obtemos conjuntos como se segue:
\begin{enumerate}
	\item[0.] $\emptyset$ %& $0$ & stage $0$\\
	\item[1.] $ \{\emptyset\}$ %& $1$ & stage $1$\\
	\item[2.] $\{\emptyset, \{\emptyset\}\}$ %& $2$ & stage $2$\\
	\item[3.] $\{\emptyset, \{\emptyset\}, \{\emptyset, \{\emptyset\}\}\}$ %&$3$ & stage $3$\\
	\item[4.] $\{\emptyset, \{\emptyset\}, \{\emptyset, \{\emptyset\}\}, \{\emptyset, \{\emptyset\}, \{\emptyset, \{\emptyset\}\}\}\}$% & $4$ & stage $4$ 
\end{enumerate}
e podemos verificar que cada um destes conjuntos é distinto. 

Começamos a numeração em $0$, por várias razões. Mas uma delas é esta. Não é
difícil verificar que o conjunto que rotulamos ``$n$'' tem exatamente $n$ membros e que (intuitivamente) se forma no $n$-ésimo estágio. 

Contudo. Isto nos dá \emph{infinitamente muitos} conjuntos, mas não
garante que exista um conjunto \emph{infinito}, isto é, um conjunto com
infinitamente muitos membros. E isto importa deveras: a menos que encontremos um
conjunto infinito de Dedekind, não podemos construir uma álgebra de Dedekind. Mas
queremos uma álgebra de Dedekind, para a podermos tratar como o conjunto dos
números naturais. (Compare-se com
\olref[sfr][infinite][dedekindsproof]{sec}.)

Importa que os axiomas que estabelecemos até aqui \emph{não}
garantem a existência de nenhum conjunto infinito. Por isso temos de enunciar um
novo axioma:

\begin{axiom}[Infinito]
Existe um conjunto, $I$, tal que $\emptyset \in I$ e $x \cup \{x\} \in I$ sempre que $x \in I$.
\begin{align*}
	\exists I( & (\exists o \in I)\forall x\ x \notin o \land {}\\
	& (\forall x \in I)(\exists s \in I)\forall z(z \in s \liff (z \in x \lor z = x)))
\end{align*}
\end{axiom}

É fácil ver que o conjunto $I$ que nos é dado pelo Axioma do
Infinito é infinito de Dedekind. O seu elemento distinguido é
$\emptyset$, e a injeção em $I$ é dada por $s(x) = x\cup
\{x\}$. Agora,
\olref[sfr][infinite][dedekind]{thm:DedekindInfiniteAlgebra} mostrou
como extrair uma álgebra de Dedekind de um conjunto infinito de Dedekind; e
trataremos isso como o nosso conjunto de números naturais. Mais precisamente:
\begin{defn}\ollabel{defnomega}
Seja $I$ qualquer conjunto que nos seja dado pelo Axioma do Infinito. Seja $s$
a função $s(x) = x \cup \{x\}$. Seja $\omega =
\closureofunder{s}{\emptyset}$. Chamamos \emph{números naturais} aos membros de $\omega$,
e dizemos que $n$ é o resultado de $n$ aplicações de $s$ a $\emptyset$.
\end{defn}

Pode agora voltar atrás e verificar que o conjunto rotulado ``$n$'', há uns
parágrafos, será tratado \emph{como} o número $n$. 

Discutiremos o significado desta estipulação na
\olref[sth][z][nat]{sec}. Por ora, permite-nos provar um
resultado intuitivo:

\begin{prop}\ollabel{naturalnumbersarentinfinite}
Nenhum número natural é infinito de Dedekind.
\end{prop}

\begin{proof}
A prova é por indução, isto é,
\olref[sfr][infinite][induction]{thm:dedinfiniteinduction}. Claramente $0
= \emptyset$ não é infinito de Dedekind. Para o passo indutivo, estabeleceremos a contrapositiva: se (absurdamente) $s(n)$ for infinito de
Dedekind, então $n$ é infinito de Dedekind. 

Suponhamos, pois, que $s(n)$ é infinito de Dedekind, isto é, que existe alguma
!!{injection} $f$ com $\ran{f}\subsetneq \dom{f} = s(n) = n \cup
\{n\}$. Há dois casos a considerar. 

\emph{Caso 1: $n \notin \ran{f}$.} Então $\ran{f} \subseteq n$, e $f(n) \in n$. Seja $g = \funrestrictionto{f}{n}$; agora $\ran{g} =
\ran{f} \setminus \{f(n)\} \subsetneq n = \dom{g}$. Logo $n$ é
infinito de Dedekind. 

\emph{Caso 2: $n \in \ran{f}$.} Fixemos $m \in \dom{f} \setminus \ran{f}$, e definamos uma função $h$ com domínio $s(n) = n \cup \{n\}$:
\[
h(x) = 
	\begin{cases}
		f(x) & \text{se }f(x) \neq n\\
		m & \text{se }f(x)=n
	\end{cases}
\]
Assim $h$ e $f$ coincidem em toda a parte, exceto que $h(f^{-1}(n)) = m
\neq n = f(f^{-1}(n))$. Uma vez que $f$ é !!a{injection}, $n \notin
\ran{h}$; e $\ran{h} \subsetneq \dom{h} = s(n)$. Logo $n$ é
infinito de Dedekind, pelo argumento do Caso~1.
\end{proof}

A questão que permanece é como poderíamos \emph{justificar} o Axioma
do Infinito. A resposta curta é que precisaremos de acrescentar outro
princípio à história que vamos contando. Esse princípio é o
seguinte:
\begin{enumerate}
	\item[] \stagesinf. Há um estágio infinito. Isto é, há um estágio que (a) não é o primeiro estágio, e que (b) tem alguns estágios antes de si, mas que (c) não tem predecessor imediato.
\end{enumerate}
O Axioma do Infinito segue-se diretamente deste princípio.
Sabemos que o número natural $n$ se forma no estágio $n$. Logo o conjunto
$\omega$ forma-se no primeiro estágio infinito. E o próprio $\omega$
testemunha o Axioma do Infinito. 

Isto, porém, apenas nos leva de volta à questão de como justificar
\stagesinf. Tal como \stagessucc, não fazia parte explícita
da história que contamos sobre a hierarquia cumulativa iterativa. Mas
mais do que isso: nada na própria ideia de uma hierarquia iterativa, em
que os conjuntos se formam estágio a estágio, nos obriga a pensar que o
processo envolve um estágio \emph{infinito}. Parece perfeitamente coerente
pensar que os estágios estão ordenados como os números naturais. 

Isto, contudo, levanta um problema óbvio. Na
\olref[sfr][infinite][dedekindsproof]{sec}, consideramos a
``prova'' de Dedekind de que existe um conjunto infinito de Dedekind (de
pensamentos). Isto poderá não lhe ter parecido muito satisfatório. Mas se
\stagesinf{} não nos é ``imposto'' pela concepção iterativa de
conjunto (nem pelas ``leis do pensamento''), continuamos sem uma
justificação intrínseca para a afirmação de que existe um conjunto infinito de
Dedekind.

Há muito mais a dizer aqui, naturalmente. Mas esperamos que esteja agora
num ponto em que possa começar a refletir sobre o que poderá \emph{custar}
justificar um axioma (ou um princípio). No que se segue tomaremos simplesmente
\stagesinf{} como concedido.

\end{document}
